\section{Related work}
\label{sec:related}

The literature was mapped against the five hypotheses \emph{before} any real-data
fit: \texttt{docs/PRIOR-ART.md} (SHA-256 \texttt{3f36fb62\ldots}), hash-stamped on
2026-08-26, forced the hypothesis-status clarifications registered that day as
Addendum~3.

Our classical arms (Supplementary Section~S2) come from a programme
\citep{lee1992modeling,girosi2007understanding,brouhns2002poisson,brouhns2005bootstrapping,li2005coherent,li2017coherent,cairns2006two,renshaw2006cohort}
that judged mortality models as density forecasters, but in-sample or ex ante, on
one or two populations, before neural models and before any break in the modern
data
\citep{cairns2009quantitative,cairns2011mortality,dowd2010evaluating,booth2008mortality};
\pkg{StMoMo} \citep{villegas2018stmomo} is our numerical oracle.

\citet{dowd2010backtesting}, this audit's direct ancestor, backtests six M1--M7
models, parameter-certain and parameter-uncertain, by locating realised mortality
inside intervals forecast from earlier origins, one population and age range at a
time. Of the two ages plotted, age~65 has ``a notable upward bias'' and age~84
``only a very slight upward bias''---better behaviour at the higher age, the
ex-ante alternative to \Hfour. A footnote grants that ``the positions of the
individual observations within the prediction intervals are not independent''
across horizons, then measures marginal positions anyway: \Hthree\ is what it
points at and does not compute. Life expectancy, survival rates and annuity prices
are named as backtestable ``in principle'' but are not backtested: \Hfive\ is
that extension.

\citet{richman2021neural} generalise Lee--Carter to all HMD populations with
country and sex embeddings, evaluated by mean squared error, as do later
architectures
\citep{hainaut2018neural,nigri2019deep,petnehazi2019mortality,perla2021time,scognamiglio2022calibrating}.
Intervals already exist in this lineage---a bootstrap ensemble of LSTMs
\citep{marino2023neural}, a variational autoencoder \citep{miyata2022extending}---so
we claim no addition of uncertainty to neural Lee--Carter.

The closest prior work is \citet{schnurch2022point}: on 54 HMD populations at ages
60--89, trained to 2006 and tested on 2007--2016, their nominal 95\% interval
coverage is 74.0 and 74.3 per cent for Lee--Carter fit on ten and twenty years,
77.2 for an autoregressive-cohort model, and 86.0, 97.0 and 98.0 for recurrent,
feed-forward and convolutional networks. Their intervals come from re-seeded
trainings of one network; the pass rule is one-sided (``minimize MPIW under the
constraint that PICP is at or above'' nominal), so over-coverage passes. The
mechanism is held fixed by design (methods needing structural change, such as
dropout layers, are ruled out), so the 74-versus-98 per cent gap is confounded between architecture and mechanism, which our crossed design
separates; their classical arms see ten or twenty training years against all
years for the networks, so part of the classical under-coverage is a
training-window artefact our expanding-origin rule removes. Their Table~2 shows
\Hone's inversion---mean squared error ranks the networks first, Poisson deviance
ranks Lee--Carter first---attributed to exposure weighting, not sharpness.

Without a break, classical intervals under-cover and bootstrap-ensemble
neural intervals over-cover \citep{schnurch2022point,marino2023neural}, so the
registered one-sided clause---``under-cover in the shift regime for every model
family''---was falsified in the stable regime before our first fit. \Htwo\ was
restated two-sided: the tested quantity is $|\text{coverage}-\text{nominal}|$,
hypothesised to grow from the stable to the shift regime, direction free by family
and mechanism, the stable regime becoming a replication arm (Addendum~3~\S8).

\citet{barigou2023bayesian} find the best model depends on country and scoring
rule, and \citet{goes2025bayesian} a top-one inversion between CRPS and the point
and log-score rankings; with \citet{schnurch2022point} they make
\Hone\ prior art, carried here as a harness-validity check, ours being only its
magnitude, per population and unweighted by exposure (Table~S2). The pandemic
literature is about \emph{estimation with} the pandemic years
\citep{cairns2020impact,robben2022assessing,goes2025bayesian,robben2024catastrophe,vanberkum2025estimating,nalmpatian2023advanced,robben2025granular},
not \emph{evaluation across} them: \citet{barigou2023bayesian} simulate a
perturbation because post-shock data did not yet exist and call a ``change of
regime'' ``beyond the scope''; \citet{goes2025bayesian} found an out-of-sample test
``impractical''. None evaluates intervals fit before 2020 against 2020--2024; with
twenty populations now final through 2024 (Table~S1) the constraint that stopped
them is gone.

Scoring and inference follow standard practice (Section~\ref{sec:metrics}), as do
the shift-adaptive and change-point conformal variants
\citep{tashman2000out,nosek2018preregistration,tibshirani2019conformal,gibbs2021adaptive,gibbs2024conformal,sun2025conformal,romano2019conformalized,angelopoulos2023gentle}.
Two things are not housekeeping: \citet{barber2023conformal} bound the
coverage gap of weighted conformal sets by a total-variation distance between the
calibration and test distributions, which a structural break is designed to make
large; and joint path coverage is solved prior art in machine learning
\citep{stankeviciute2021conformal,sun2024copula,messoudi2021copula,diquigiovanni2022conformal}.

Conformal prediction reached mortality forecasting only in a 2025--2026 line
\citep{shang2025intervals,shang2026conformal,shang2026subnational}: all three use
one functional forecasting family and evaluate coverage in ordinary years; none
audits conformal coverage across a designated break, places conformal beside other
uncertainty mechanisms, or measures joint path coverage, and in our 52-paper corpus
``conformal'', ``Monte Carlo dropout'' and ``deep ensemble'' have no hit. Three conformal mechanisms therefore enter our grid as
\emph{audited arms}; we claim the first \emph{measurement} of the marginal-to-joint
coverage gap in mortality forecasting, not the method, make no claim that any
conformal arm fixes anything, and do not invoke its finite-sample guarantee under a
break.

What remains is the crossed family-by-mechanism design, the conformal arms across
the break, the marginal-to-joint gap, the derived-quantity coverage audit, and the
pandemic test the field declined for lack of data. Two items were unread, named
so their absence is visible:
\citet{schnurch2022point}'s online supplement (Section~C.3, implied annuity present
values), the one live threat to \Hfive's residual novelty should it audit
annuity-interval coverage rather than report point values; and the full text of
\citet{stankeviciute2021conformal}, cited from later work.
